\chapter{The critical strip and Riesz's criterion}\label{chap:strip}

Chapter~\ref{chap:basepoint} read the reciprocal Xi function at a basepoint $\alpha$ as a Laplace transform on the half line, $\varphi_\alpha(u)=\xi(\alpha)/\xi(\alpha+\sqrt u)$ for $u\ge0$. At the basepoint one its Thorin measure is the Poisson smoothing of the zeros and its L\'evy form carries the primes, and by Theorem~\ref{bp:thm:alpha} it is, for every $\alpha>\frac12$ and with no hypothesis, the Laplace transform of a positive random variable $T_\alpha$ obtained by subordinating the centre clock. The same function has two further readings, as the characteristic function of a normal scale mixture and through the completeness of the translates of the tilted P\'olya kernels. Both run across the whole critical strip, not only to the right of the zeros, and so show how much of what holds with no hypothesis at the basepoint one survives inside the strip. Section~\ref{ch:strip} develops the two readings in turn and ends with Riesz's criterion, which reads the zero free half planes through growth rather than completeness. Section~\ref{ch:riesz} writes that criterion in the language of the clocks, where it becomes the statement that every $T_\alpha$ with $\alpha>\frac12$ is a GGC, and shows why none of the resulting representations of its coefficients supplies the growth bound that the criterion needs. Section~\ref{sec:sqfree} turns to the squarefree integers, whose Dirichlet series $\zeta(s)/\zeta(2s)$ carries the same reciprocal $1/\zeta(2s)$: the constant of $\sum_{n\le x}\mu(n)^2/n$ is a combination of two tilted P\'olya means, and smoothing the count by a gamma law gives criteria for RH of the same kind as Riesz's.

For the scale mixtures and the completeness statements, as for the Thorin measure and the L\'evy form in Chapter~\ref{chap:basepoint}, what is unconditional at the basepoint one extends to every basepoint to the right of the line $\re s=\alpha$, for a given $\alpha<1$, exactly when no zero lies to the right of that line; required for every $\alpha>\frac12$, this is RH. Chapters~\ref{chap:basepoint} to \ref{chap:almost} prove no new zero free region.
